\chapter{The Rational Numbers}

Over the set \(\Z \times (\Z \setminus \{0\})\) we define the relation

\[
    (a, b) \sim (c, d) \iff ad = bc
\]

which is an equivalence relation. The set of all equivalence classes is called \(\Q\) as we write 

\[
    \frac{a}{b} := [(a, b)].
\]

The choice of the raltion may seem arbitrary, but if we return to our intuitive defintion for rational number where 

\[
    \frac{a}{b} = \frac{c}{d} \iff c = af\, \text{ and } \, d = bf
\]

we see that that iff we divide both fractions we get

\[
    \frac{ad}{bc} = \frac{abf}{abf} = 1
\]

Meaning that we were dividing by our same number.

\section{Definition of Addition and Multiplication}

For \(a, c \in \Z\) and \(b, d \in \Z \setminus \{0\}\) we define 

\[
    \frac{a}{b} + \frac{c}{d} := \frac{ad + bc}{bd}, \quad \frac{a}{b} \cdot \frac{c}{d} := \frac{ac}{bd}
\]

Both operations are independent from the choice of the presentants of the class.